\section{Recommender Systems}

\subsection{Collaborative Filtering}

Collaborative filtering (CF) is one of the most established paradigms in recommender systems. Its central assumption is that user preference can be inferred from historical interactions: users who behaved similarly in the past are likely to prefer similar items in the future, and items consumed by similar groups of users are likely to be related. Under this view, recommendation is commonly formulated as a preference prediction or ranking problem. Given a set of users \(\mathcal{U}\), a set of items \(\mathcal{I}\), and an observed interaction matrix \(R \in \mathbb{R}^{|\mathcal{U}| \times |\mathcal{I}|}\), the objective is to learn a scoring function
\begin{equation}
    f(u,i) \rightarrow \hat{r}_{ui},
\end{equation}
where \(\hat{r}_{ui}\) denotes the predicted relevance of item \(i \in \mathcal{I}\) to user \(u \in \mathcal{U}\). The learned score can then be used to estimate missing entries in \(R\) or to generate a top-\(N\) ranked list of items for each user. Recent surveys emphasize that this formulation remains the foundation of many modern graph-based CF models, although the representation and propagation mechanisms have become increasingly sophisticated \parencite{11359174}.

The interaction matrix \(R\) can be constructed from either explicit or implicit feedback. Explicit feedback directly reflects user opinions, such as numerical ratings, likes, or preference labels. In this setting, each observed value \(r_{ui}\) usually indicates the degree of satisfaction of user \(u\) with item \(i\). In contrast, implicit feedback is derived from behavioral signals, such as clicks, views, purchases, add-to-cart events, or check-ins. These signals are easier to collect at scale but are more ambiguous, because an observed interaction does not necessarily imply strong preference and an unobserved entry does not necessarily indicate dislike. Consequently, implicit-feedback recommendation is often treated as a ranking task rather than a rating-regression task. This distinction is particularly important in e-commerce and fashion recommendation, where user decisions are frequently represented by sparse behavioral logs rather than explicit ratings.

Traditional CF methods can be broadly divided into memory-based and model-based approaches. Memory-based CF operates directly on the observed interaction matrix and produces recommendations by comparing users or items according to their historical interaction patterns. In user-based CF, the preference of a target user is predicted from the preferences of other users with similar interaction histories. The underlying intuition is that users with similar past behaviors tend to share future interests. Conversely, item-based CF estimates the relevance of a candidate item by identifying items similar to those that the target user has previously interacted with. This strategy is often more stable in practical systems when the item catalog changes more slowly than user behavior.

The effectiveness of memory-based CF depends strongly on the similarity measure used to compare users or items. Common measures include cosine similarity, Pearson correlation, and Jaccard similarity. Cosine similarity measures the angle between two interaction vectors and is widely used when the magnitude of interactions is less important than their direction. Pearson correlation removes the effect of different rating scales by centering ratings around user or item means. Jaccard similarity, in contrast, is typically applied to binary implicit-feedback data and measures the overlap between two sets of interacted items or users. Although these measures are conceptually simple and interpretable, they rely heavily on sufficient co-occurrence evidence. As a result, their performance often degrades when the interaction matrix is highly sparse.

Model-based CF addresses this limitation by learning compact latent representations of users and items from the observed interactions. Matrix factorization (MF) is a representative model-based technique. It approximates the interaction matrix as the product of two low-dimensional embedding matrices,
\begin{equation}
    R \approx U V^{T},
\end{equation}
where \(U \in \mathbb{R}^{|\mathcal{U}| \times d}\) contains user latent factors, \(V \in \mathbb{R}^{|\mathcal{I}| \times d}\) contains item latent factors, and \(d\) is the embedding dimension. The predicted relevance score is commonly computed by the inner product \(\hat{r}_{ui} = \mathbf{u}_{u}^{T}\mathbf{v}_{i}\), where \(\mathbf{u}_{u}\) and \(\mathbf{v}_{i}\) are the latent vectors of user \(u\) and item \(i\), respectively. Compared with memory-based methods, MF can generalize beyond directly observed co-occurrences because users and items are projected into a shared latent space. This property makes MF more robust to moderate sparsity and more scalable for large recommendation scenarios.

Despite its effectiveness, conventional MF still represents user--item interactions primarily through pairwise latent factors. It does not explicitly model the higher-order collaborative structure formed by multi-hop relations among users and items. For example, two users may not share many directly interacted items, but they can still be related through a chain of intermediate users and items. Similarly, two items may be semantically or behaviorally connected even if they have limited direct co-consumption. Such high-order connectivity is naturally represented as a bipartite graph, where users and items are nodes and observed interactions are edges. This observation has motivated the integration of graph neural networks (GNNs) into CF, allowing collaborative signals to be propagated over the user--item graph \parencite{11359174}.

LightGCN is a representative example of this direction. It simplifies graph convolution for recommendation by removing feature transformation and nonlinear activation, and focuses on neighborhood aggregation over the user--item interaction graph \parencite{he2020lightgcnsimplifyingpoweringgraph}. From a CF perspective, LightGCN can be interpreted as extending latent factor models with graph-based embedding propagation. Instead of learning user and item embeddings only from direct interactions, it refines them by aggregating information from multi-hop neighbors. This design highlights an important transition in recommender-system research: from static latent-factor matching to structural representation learning over interaction graphs.

Nevertheless, both memory-based and classical model-based CF suffer from several well-known limitations. First, data sparsity remains a fundamental challenge because most users interact with only a small fraction of available items. Sparse observations reduce the reliability of similarity estimation and limit the quality of learned embeddings. Second, the cold-start problem arises when new users or new items have few or no interactions. Since CF primarily depends on historical behavior, it cannot effectively represent such users or items without additional information. Third, traditional CF has limited ability to exploit heterogeneous side information, such as product images, textual descriptions, categories, prices, or social relations. This limitation is especially restrictive in fashion recommendation, where visual appearance, style compatibility, and semantic descriptions are central to user preference.

These limitations provide the technical motivation for graph-based and multimodal recommender systems. Graph-based CF improves the modeling of high-order collaborative signals by treating the interaction matrix as a graph structure, while multimodal recommendation incorporates auxiliary content features to alleviate sparsity and cold-start issues. In the context of this study, CF serves as the methodological foundation: it defines the user--item preference prediction problem, explains the role of interaction data, and motivates the use of LightGCN and multimodal item representations in later sections.
